\section{Exact Gamma and Bell formulas for the Laurent coefficients}
\label{clc:sec:coefficients}

The finite closure becomes an explicit algorithm once the coefficients
$a_{\rho,k}$ are known. For a fixed pole put
\[
 I_\rho=\{j:\rho/q_j\in\mathbb Z\},\qquad
 J_\rho=\{1,\ldots,r\}\setminus I_\rho,
 \qquad d=|I_\rho|.
\]
Define the complete exponential Bell polynomials by
\begin{equation}\label{clc:eq:Bell-def}
 \exp\left(\sum_{n\ge1}u_n\frac{z^n}{n!}\right)
 =\sum_{n\ge0}\mathcal B_n(u_1,\ldots,u_n)\frac{z^n}{n!}.
\end{equation}
Thus $\mathcal B_0=1$, $\mathcal B_1=u_1$,
$\mathcal B_2=u_1^2+u_2$, and
$\mathcal B_3=u_1^3+3u_1u_2+u_3$.

\begin{theorem}[Finite coefficient construction]\label{clc:thm:coeff}
The leading coefficient at $\rho$ is
\begin{equation}\label{clc:eq:leading}
 a_{\rho,d}=(-1)^{\sum_{j\in I_\rho}\rho/q_j}
 \prod_{j\in J_\rho}\frac{\pi}{q_j}\csc\frac{\pi\rho}{q_j}.
\end{equation}
For $n\ge1$ let
\begin{align}
 u_n={}&\boldsymbol1_{2\mid n}\,2(n-1)!\zeta(n)
              \sum_{j\in I_\rho}q_j^{-n}\notag\\
 &+\sum_{j\in J_\rho}q_j^{-n}
 \left\{\psi^{(n-1)}(\rho/q_j)
       +(-1)^n\psi^{(n-1)}(1-\rho/q_j)\right\},
 \label{clc:eq:logjets}
\end{align}
where the first line is zero for odd $n$, so it never calls for
$\zeta(1)$. Then, for $1\le k\le d$,
\begin{equation}\label{clc:eq:Bell-coeff}
 a_{\rho,k}=\frac{a_{\rho,d}}{(d-k)!}
                  \mathcal B_{d-k}(u_1,\ldots,u_{d-k}).
\end{equation}
Only $u_1,\ldots,u_{d-1}$ are needed. For rational scales,
$a_{\rho,k}/\pi^{r-k}$ is a real cyclotomic algebraic number.
\end{theorem}
\begin{proof}
Each polar factor has leading term
$(-1)^{\rho/q_j}/z$, proving \eqref{clc:eq:leading}.
The normalized local germ
$z^dG_{\bq}(\rho+z)/a_{\rho,d}$ is analytic and equals one at zero;
its local logarithm is consequently unambiguous.
For a polar factor its logarithm is
\[
 \log\frac{\pi z/q_j}{\sin(\pi z/q_j)}
 =\sum_{h=1}^\infty\frac{\zeta(2h)}{h}\left(\frac z{q_j}\right)^{2h}.
\]
This follows from the sine product, or from Gamma reflection and the
local Gamma series. Its $n$th derivative at zero is the first line of
\eqref{clc:eq:logjets}.
For a nonpolar factor use
\[
 g_{q_j}(\rho+z)=q_j^{-1}
 \Gamma\left(\frac{\rho+z}{q_j}\right)
 \Gamma\left(1-\frac{\rho+z}{q_j}\right).
\]
Differentiating its local logarithm gives the second line. The arguments
$1-\rho/q_j$ here may be negative, but are not integers; the meromorphic
polygamma values are well defined. No global branch of their separate
Gamma logarithms is required. Exponentiating the normalized local
logarithm gives \eqref{clc:eq:Bell-coeff}.

For the final assertion one can avoid transcendental cancellation
entirely. Expand the sine factors directly at rational multiples of
$\pi$, invert their finite Taylor polynomials, and multiply to order
$d-1$. The trigonometric constants are real cyclotomic algebraic
numbers, and the Taylor powers supply exactly $\pi^{r-k}$.
\end{proof}

The two constructions in this proof are independent computational
routes. The supplied code uses the logarithmic Bell recurrence in
\texttt{pole\_table}, and direct reciprocal-sine Taylor multiplication in
\texttt{direct\_local\_table}. Exact comparisons are made after symbolic
simplification, not after floating-point evaluation of polygamma terms.

\subsection{A genuinely mixed three-scale example}
For $\bq=(1,2,3)$, a common period is $Q=6$ and $\sigma=-1$.
The complete table is as follows; blank entries mean zero or absent
higher pole order.
\begin{center}
\begin{tabular}{c c c c}
\toprule
$\rho$ & $a_{\rho,1}$ & $a_{\rho,2}$ & $a_{\rho,3}$\\
\midrule
$1$ & $-\sqrt3\pi^2/9$ & &\\
$2$ & $-2\pi^2/27$ & $-2\sqrt3\pi/9$ &\\
$3$ & $0$ & $-\pi/2$ &\\
$4$ & $2\pi^2/27$ & $-2\sqrt3\pi/9$ &\\
$5$ & $\sqrt3\pi^2/9$ & &\\
$6$ & $-49\pi^2/216$ & $0$ & $-1$\\
\bottomrule
\end{tabular}
\end{center}
Together with \eqref{clc:eq:central-raw}, this is an explicit finite
identity for every total order and every $\Rea A>-1$, using alternating
Hurwitz functions at $(A+\rho)/6$. No additional infinite sum is hidden
in the coefficient construction.
At $W=0$ every $k\ge2$ term vanishes because $(0)_{k-1}=0$, and
$\eta(0,b)=1/2$. Thus
\begin{equation}\label{clc:eq:123zero}
 K_{(1,2,3)}(0)=-\frac12\sum_\rho a_{\rho,1}
 =\frac{49\pi^2}{432}.
\end{equation}
The independent parity theorem in Section~\ref{clc:sec:parity} gives
exactly the same value without constructing this table.

The coefficients also obey a useful scaling check. If all scales and
poles are multiplied by $\lambda>0$, then
\begin{equation}\label{clc:eq:coef-scaling}
 a^{\lambda\bq}_{\lambda\rho,k}
 =\lambda^{k-r}a^{\bq}_{\rho,k}.
\end{equation}
This follows at once from
$G_{\lambda\bq}(p)=\lambda^{-r}G_{\bq}(p/\lambda)$ and is included
in the exact verification suite.
